\section{Basics about collar neighbourhoods}
\label{sec:mathsdef}

\begin{defi}
    A \emph{topological manifold} is a second-countable Hausdorff space that is locally Euclidean.
\end{defi}

\begin{rem}
    Apparently, there are some subtleties with collars of non-metrizable manifolds but I assume we don't care. (At least with the above definition our manifolds are always metrizable.)
\end{rem}

\begin{defi}
    Let $M$ and $N$ be two $C^k$-manifolds. 
    Let $f \colon M \to N$. 
    \begin{enumerate}
        \item $f$ is a $C^k$-diffeomorphism if it is bijective and both $f$ and its inverse functions are $k$-times continuously differentiable. 
        \item $f$ is a $C^k$-embedding if it is a $C^k$-diffeomorphism onto its image.
    \end{enumerate}
\end{defi}

\begin{defi}[Due to \cite{Connelly1971}]\label{defi:collared}
    Let $M$ be a $\fC^k$-manifold (with boundary) and let $B \subseteq M$ be closed. 
    \begin{enumerate}\label{defi:collared:local}
        \item $B$ is \emph{locally collared} if it is covered by sets $U \subseteq B$ which are open in $B$ such that there are $\fC^k$-embeddings $\varphi_U \colon \closure{U} \times [0, 1] \to M$ with 
        \begin{enumerate}
            \item $\preim{\varphi_U}{B} = \closure{U} \times \{ 0\}$.
            \item For all $x \in \closure{U}$, $\varphi_U (x, 0) = x$.
            \item $\varphi_U (\closure{U} \times [0, 1])$ is closed. 
            \item $\varphi_U (U \times [0, 1)$ is open.  
        \end{enumerate}
        \item $B$ is \emph{collared} if $B$ already fulfils the conditions given above. In this case we call the image of $B \times [0, 1)$ under the corresponding map $\varphi_B \colon B \times [0, 1] \to M$ a \emph{collar neighbourhood} of $B$.
    \end{enumerate}
\end{defi}

\begin{rem}
    This definition is a lot more restrictive than other definition which don't consider the closure or the closed interval. I think if this one doesn't work, I can just take the other one. 
\end{rem}

\begin{rem}
    For $k=0$ we can even extend this definition to topological spaces instead of just defining it for topological manifolds. 
\end{rem}

\begin{thm}[Due to \cite{Connelly1971}]
    Let $X$ be a Hausdorff space and $B \subseteq X$ be a closed subset. 
    If $B$ is locally collared in $X$, then it is collared.
\end{thm}

\begin{proof}
    Let $U_1, \dots , U_n$ be a cover of $B$ that witnesses that $B$ is locally collared. 
    Let $\varphi_i \colon \closure{U_i} \times [0, 1] \to X$ be the corresponding embeddings. 
    $B$ is normal because it is compact and Hausdorff. 
    By the shrinking lemma (already in mathlib) there is a new open cover $V_1, \dots, V_n$ of $B$ such that $\closure{V_i} \subseteq U_i$ for all $1 \leq i \leq n$. 
    Define 
    \begin{equation*}
        X^+ \coloneq (X \amalg B \times [-1, 0]) / x \sim (x, 0).
    \end{equation*}
    Clearly, $X^+$ is collared. So it suffices to show that there is a homeomorphism $G \colon X \to X^+$ s.t.\ $G(x) = (x, -1)$ for all $x \in B$. 
    Let $\Phi_i \colon \closure{U_i} \times [-1, 1] \to X^+$ be defined by 
    \begin{equation*}
        \Phi_i (x) = \begin{cases*}
            \phi_i(x) & \text{if } $x \in \closure{U_i} \times [0, 1]$ \\
            x & \text{if } $x \in \closure{U_i} \times [-1, 0]$. \\
        \end{cases*}
    \end{equation*}
    We now inductively define maps $f_i \colon B \to [0, 1]$ and embeddings $g_i \colon X \to X^+$ for all $0 \leq i \leq n$ such that 
    \begin{enumerate}
        \item $f_i(x) = -1$ if $x \in \bigcup_{j \leq i} \closure{V_j}$
        \item $g_i(x) = (x, f_i(x))$ if $x \in B$
        \item $g_i(X) = X \cup \{ (x, t) \mid t \geq f_i(x)\}$
    \end{enumerate}
    Then $g_n$ will be the desired homeomorphism $G$ from above.
    Set $f_0$ to be constant $0$. Then we can choose $g_0$ to be the inclusion and trivially $g_0(X) = X = X \cup X \times \{0\} = X \cup \{ (x, t) \mid t \geq 0\}$. 
    
    Assume that $g_{i-1}$ has been defined. Using Urysohn's Lemma, choose a continuous function $\lambda_i \colon \closure{U_i} \to [0, 1]$ such that it is $0$ on $\closure{U_i} \setminus U_i$ and $1$ on $\closure{V_i}$. 
    Given $x \in B$ set $a(x) \coloneq f_{i-1}(x)$ and $b(x) \coloneq (1 - \lambda_i(x))f_{i-1}(x) + \lambda_{i}(x) (-1)$.
    Now let 
    \begin{equation*}
        s(x, \_) \colon [a(x), 1] \to [b(x), 1],\quad t \mapsto \frac{b(x)-1}{a(x)-1}(t-1) + 1.
    \end{equation*}
    Now set
    \begin{equation*}
        \psi_i \colon \preim{\Phi_i}{g_{i-1}(X)} \to \closure{U_i} \times [-1,1], \quad (x, t) \mapsto (x, s(x, t)).
    \end{equation*}
    This allows us to define 
    \begin{equation*}
        \Psi_i \colon g_{i-1}(X) \to X^+, \quad x \mapsto \begin{cases*}
            \Phi_i (\psi_i (\inv{\Phi_i}(x))) & if  $x \in g_i(X) \cap \Phi_i(\closure{U_i} \times [-1, 1])$ \\ 
            x & \text{otherwise} \\
        \end{cases*}
    \end{equation*}
    \begin{claim}
        $\Psi_i$ is continuous. 
    \end{claim}
    \begin{proof}
        It suffices to show that $\psi_i$ is continuous. 
        We need to show that $s \colon B \times [a(x), 1] \to [b(x), 1]$ is continuous. 

        Ahhh this is ugly. I somehow would prefer to just get going with this in Lean.
    \end{proof}
\end{proof}

I want to try to separate out some parts here: 

\begin{defi}
    Let $X$ be a $\fC^k$-manifold and $B \subseteq X$ a closed subset. 
    We define the topological space
    \begin{equation*}
        X^+_B \coloneq (X \amalg B \times [-1, 0]) / x \sim (x, 0).
    \end{equation*}
\end{defi}

\begin{lem}
    Let $X$ be an $n$-dimensional $\fC^k$-manifold and $B \subseteq X$ a closed subset. 
    Then $X^+_B$ is an $n$-dimensional $\fC^k$-manifold. 
\end{lem}
\begin{proof}
    We first need to show that $X^+_B$ is a Hausdorff space. I am not sure if this is best done manually or if there is some cool lemma like: a pushout along closed embeddings preserves Hausdorff spaces. (Is this true?)

    Then we need to show that $X^+_B$ is second countable. Again, there should be some general lemma about closed embeddings and pushouts. 

    Now we need to pick an atlas for $X^+_B$.
    Let $\fA_X$ be a $\fC^k$-atlas for $X$, $\fA_B$ be a $\fC^k$-atlas for $B$ and $\fA_{[-1, 0]}$ be a $\fC^k$-atlas for $[-1, 0]$. 
    We define 
    \begin{align*}
        \fA \coloneq &\{(U, \varphi) \in \fA_X \mid U \cap B = \varnothing\} \\
        &\cup \{(U \times I, \varphi \times \psi) \mid (U, \varphi) \in \fA_B, (I, \psi) \in \fA_{[-1, 0]}, 0 \notin I\} \\
        &\cup \{(U \cup ((U \cap B) \times I), \varphi^+_I) \mid (U, \varphi) \in \fA, (I, \psi) \in \fA_{[-1, 0]}, 0 \in I\}
    \end{align*}
    where we define $\varphi^+_I$ as follows. 
    Let $(U, \varphi)$ and $(I, \psi)$ as above. 
    Set 
    \begin{equation*}
        \varphi^+_I \colon U \cup ((U \cap B) \times I) \to \bR^n, x \mapsto \begin{cases*}
            \varphi (x) & $x \in U$ \\ 
            (\varphi(y), z) & $x = (y, z) \in (U \cap B) \times I$ \\
        \end{cases*}
    \end{equation*}
    For $x \in B$ we have $\varphi(x)^+_I = \varphi(x) = (\varphi(x), 0) = \varphi^+_I(x, 0)$, so $\varphi^+_I$ is continuous. 

    \textbf{It seems to me that this might not be $\fC^k$ at all. I don't know of you can even glue this without basically already having the vector field available and then one might as well do the other proof. Maybe a local vector field is enough? Not sure :(}

    Now you would need to verify that this is a $\fC^k$ atlas. 
    Possibly using a version of the manifold chart lemma from Lee. 
    I will only try to show that the transition maps are $\fC^k$. 
    The rest should be easier. 
    The only interesting intersections are those of two charts of the form $(U \cup ((U \cap B) \times I), \varphi^+_I)$ and $(V, \cup ((V \cap B) \times J), \psi^+_J)$. 
    We need to check that the transition map is $\fC^k$ on $U \cap V \cap B$. 
    This is clear for all combinations of partial derivatives that do not include the $n$-th one. 
\end{proof}

\begin{lem}
    Let $X$ be an $n$-dimensional $\fC^k$-manifold and $B \subseteq X$ a closed subset. 
    Then $X^+_B$ is the pushout in the category of $\fC^k$-manifolds. 
\end{lem}

\begin{rem}
    It seems that this is not at all trivial and in fact much of the issue with gluing manually along boundaries.
    I am not sure if we can even guarantee the above lemma.
    The above lemma would hopefully mean that one never has to look at the explicit construction again.
\end{rem}


\begin{defi}
    Let $X$ be a Hausdorff space and $B \subseteq X$ a closed subset.
    Let $U \subseteq B$ be an open subset together with a 
\end{defi}

We would like the following theorem: 

\begin{thm}
    Let $M$ be a $\fC^k$-manifold (with boundary) and $B \subseteq X$ a closed subset. 
    If $B$ is locally collared in $M$, then it is collared.
\end{thm}

Then we will need the following lemma: 

\begin{lem}
    Let $M$ be an $n$-dimensional $\fC^k$-manifold with boundary. 
    Then $\boundary M$ is locally collared.
\end{lem}

\begin{proof}
    Let $p \in \boundary M$. Then $p$ is contained in a regular half-coordinate ball $U$, i.e.\ there is a coordinate half-ball $U' \supseteq \closure{U}$ and a coordinate map $\psi \colon U' \to \bH^n$ such that for some $r' > r > 0$ we have $\psi(U) = B_r(0) \cap \bH^n$, $\psi (\closure{U}) = \closure{B_r(0)} \cap \bH^n$ and $\psi (U') = B_{r'}(0) \cap \bH^n$. 

    Now pick an open neighbourhood $V$ of $0 \in \boundary \bH^n$ and $l > 0$ s.t. $V \times [0, l] \subseteq B_r(0)$.

    Considering $\closure{V}$ and $V$ as submanifolds of $\bH^n$ and $\preim{\psi}{\closure{V}} = \closure{\preim{\psi}{V}}$ and $\preim{\psi}{V}$ as submanifolds. 
    Then 
    \begin{equation*}
        \restrict{\psi}{\closure{\preim{\psi}{V}}} \colon \closure{\preim{\psi}{V}} \to \closure{V}
    \end{equation*}
    and 
    \begin{equation*}
        \restrict{\psi}{\preim{\psi}{V}} \colon \preim{\psi}{V} \to V
    \end{equation*}
    are $\fC^k$-diffeomorphisms. 
    So it suffices to show that $\restrict{\inv{\psi}}{\closure{V} \times [0, l]} \colon \closure{V} \times [0, l] \to M$ fulfils the conditions given in Definition \ref{defi:collared}. 

    I think this should work... 
    I think it would be nicer to do this directly in Lean and then later write up the maths proof. 
\end{proof}

which will then allow us to derive

\begin{thm}[Collar neighbourhood theorem]
    Let $M$ be a $C^k$-manifold with boundary, then $\boundary M$ is collared. 
\end{thm}

%\begin{defi}
%    Let $M$ be a $\fC^k$-manifold with boundary and let $U$ be a neighbourhood of the boundary $\boundary M$. 
%    $U$ is called a \emph{collar neighbourhood} if there is a $\fC^k$-diffeomorphism $f \colon \boundary M \times [0, 1) \to U$ s.t.\ for all $x \in \boundary M$, $f(x, 0) = x$. 
%\end{defi}